\documentclass[12pt,a4paper]{article}

\usepackage[utf8]{inputenc}
\usepackage[T2A]{fontenc}
\usepackage[russian]{babel}
\usepackage{amsmath, amssymb}
\usepackage{graphicx}
\usepackage{booktabs}
\usepackage{array}
\usepackage{geometry}
\usepackage{hyperref}
\usepackage{float}
\usepackage{caption}
\usepackage{subcaption}
\usepackage{siunitx}
\usepackage{indentfirst}
\usepackage{pgfplots}

\geometry{
  left=2.5cm, right=2cm,
  top=2cm, bottom=2cm
}

\sisetup{
  output-decimal-marker={,},
  group-separator={~},
  group-minimum-digits=4
}

\title{\textbf{Отчёт по лабораторной работе}\\
\large Низкочастотные процессы в многомодовом твердотельном лазере}

\author{Актанаев Александр \and Кукушкин Юрий \and Можаров Андрей}
\date{\today}

\begin{document}

\maketitle
\tableofcontents
\newpage

%%%%%%%%%%%%%%%%%%%%%%%%%%%%%%%%%%%%%%%%%%%%%
\section{Цель работы}

Практическое ознакомление с основными характеристиками излучения непрерывного
многомодового лазера на кристалле алюмоиттриевого граната (Nd:YAG), накачиваемого
излучением полупроводникового лазерного диода. Исследование флуктуаций полной
интенсивности излучения и поведения синфазных релаксационных колебаний.

%%%%%%%%%%%%%%%%%%%%%%%%%%%%%%%%%%%%%%%%%%%%%
\section{Экспериментальной установки}

В качестве источника накачки используется полупроводниковый лазер GaAs/GaAlAs со следующими характеристиками:
длина волны генерации $\lambda = 810$~нм,
пороговый ток питания $I_\text{пор} = 200$~мА,
максимальный ток питания $I_{max} = 800$~мА,
максимальная мощность излучения $P_{max} = 0{,}5$~Вт.

Активной средой лазера является кристалл YAG, легированный ионами
$\mathrm{Nd}^{3+}$ с концентрацией 1\%. Кристалл имеет форму цилиндра длиной
$l = 1$~мм. Коэффициент отражения выходного зеркала
$R_2 = 0{,}98\ldots0{,}995$ на длине волны $\lambda = 1064$~нм.

%%%%%%%%%%%%%%%%%%%%%%%%%%%%%%%%%%%%%%%%%%%%%
\section{Основные элементы теории}

\subsection{Балансные уравнения многомодового лазера}

Динамика многомодового Nd:YAG лазера описывается балансными уравнениями
Танга, Статца и Де Марса (ТСД)~\cite{tang1963}:
\begin{align}
  \frac{dI_p}{d\tau} &= G I_p \bigl[\gamma_p(N_0 - N_p) - 1\bigr], \label{eq:I}\\
  \frac{dN_0}{d\tau} &= A - \Bigl(1 + \sum_{p=1}^{M}\gamma_p I_p\Bigr)N_0
                        + \sum_{p=1}^{M}\gamma_p I_p N_p, \label{eq:N0}\\
  \frac{dN_p}{d\tau} &= -\Bigl(1 + \sum_{p=1}^{M}\gamma_p I_p\Bigr)N_p
                        + \frac{1}{2}\gamma_p I_p N_0. \label{eq:Np}
\end{align}
Здесь $I_p$~--- интенсивность $p$-й моды; $\tau = t/T_1$~--- безразмерное
время; $G = T_1/T_c$~--- отношение времени жизни верхнего уровня $T_1 \approx 0{,}23$~мс к
времени жизни фотона в резонаторе $T_c$; $A = P/P_\text{пор}$~---
параметр накачки; $N_0$, $N_p$~--- однородная компонента и $p$-я гармоника
инверсии населённостей.

\subsection{Стационарная генерация}

Нетривиальное стационарное решение при $A > 1$ (условие самовозбуждения):
\begin{equation}
  \bar{I} = A - 1, \qquad \bar{N}_0 = 1.
  \label{eq:stationary}
\end{equation}
Это означает, что мощность лазерного излучения линейно растёт с превышением
накачки над порогом.

\subsection{Релаксационные колебания}

Линеаризация уравнений вблизи стационарного состояния даёт затухающие
колебания~--- релаксационные колебания~--- с безразмерной частотой:
\begin{equation}
  \Omega_R = \sqrt{G(A-1)},
  \label{eq:OmegaR}
\end{equation}
и декрементом затухания $\Theta_R = -A/2$.

Реальная (размерная) частота релаксационных колебаний:
\begin{equation}
  f_R = \frac{\Omega_R}{2\pi T_1} = \frac{1}{2\pi T_1}\sqrt{G(A-1)}.
  \label{eq:fR}
\end{equation}

\subsection{Полоса резонатора}

Из параметра $G$ определяется время жизни фотона $T_c = T_1/G$ и полоса
(ширина линии) резонатора:
\begin{equation}
  \delta f_\text{эксп} = \frac{1}{2\pi T_c} = \frac{G}{2\pi T_1}.
  \label{eq:deltaf}
\end{equation}

Теоретическое значение полосы резонатора определяется потерями на зеркалах:
\begin{equation}
  \delta f_\text{теор} = -\frac{c\ln\sqrt{R_1 R_2}}{2\pi L},
  \label{eq:deltaf_teor}
\end{equation}
где $c$~--- скорость света, $L$~--- длина резонатора, $R_1$, $R_2$~---
коэффициенты отражения зеркал.

\subsection{Параметр накачки}

Для лазера Nd:YAG с полупроводниковым источником накачки:
\begin{equation}
  A_{\rm Nd:YAG} = \frac{P}{P_\text{пор}},
  \label{eq:A}
\end{equation}
где $P_\text{пор}$~--- пороговая мощность накачки.

%%%%%%%%%%%%%%%%%%%%%%%%%%%%%%%%%%%%%%%%%%%%%
\section{Параметры установки}

\begin{table}[H]
  \centering
  \caption{Параметры лазерной системы}
  \label{tab:params}
  \begin{tabular}{lll}
    \toprule
    Параметр & Обозначение & Значение \\
    \midrule
    Пороговая мощность накачки & $P_\text{пор}$ & $164$~мВт \\
    Пороговая (опорная) частота & $f_\text{пор}$ & $289$~кГц \\
    Длина кристалла & $l$ & $1$~мм \\
    Время жизни верхнего уровня Nd:YAG & $T_1$ & $0{,}23$~мс \\
    Длина волны генерации & $\lambda$ & $1064$~нм \\
    Длина резонатора & $L$ & $\sim 6$~см \\
    Коэффициент отражения выходного зеркала & $R_2$ & $0{,}99$ \\
    \bottomrule
  \end{tabular}
\end{table}

%%%%%%%%%%%%%%%%%%%%%%%%%%%%%%%%%%%%%%%%%%%%%
\section{Результаты эксперимента}

\subsection{Задание~1: Зависимость релаксационной частоты от параметра накачки}

Снята зависимость релаксационной частоты $f_R$ от мощности накачки $P$
(не менее~12 точек). Параметр накачки вычислен по формуле~\eqref{eq:A}.
Для каждой точки из формулы~\eqref{eq:fR} вычислен параметр $G$:
\begin{equation}
  G_i = \frac{\bigl(2\pi f_{R,i} T_1\bigr)^2}{A_i - 1}.
  \label{eq:G_calc}
\end{equation}

Результаты измерений и вычислений представлены в таблице~\ref{tab:task1}.

\begin{table}[H]
  \centering
  \caption{Зависимость релаксационной частоты от параметра накачки и расчёт $G$}
  \label{tab:task1}
  \begin{tabular}{cccccc}
    \toprule
    № & $P$, мВт & $f_R$, кГц & $A = P/P_\text{пор}$ & $A-1$ & $G$\\
    \midrule
    1  & 180 & 445,3 & 1,0976 & 0,0976 & $4{,}24\times10^6$ \\
    2  & 187 & 500,0 & 1,1402 & 0,1402 & $3{,}72\times10^6$ \\
    3  & 196 & 550,7 & 1,1951 & 0,1951 & $3{,}25\times10^6$ \\
    4  & 214 & 613,2 & 1,3049 & 0,3049 & $2{,}58\times10^6$ \\
    5  & 233 & 648,4 & 1,4207 & 0,4207 & $2{,}09\times10^6$ \\
    6  & 254 & 683,5 & 1,5488 & 0,5488 & $1{,}78\times10^6$ \\
    7  & 272 & 699,2 & 1,6585 & 0,6585 & $1{,}55\times10^6$ \\
    8  & 287 & 738,2 & 1,7500 & 0,7500 & $1{,}52\times10^6$ \\
    9  & 307 & 765,6 & 1,8720 & 0,8720 & $1{,}40\times10^6$ \\
    10 & 317 & 796,8 & 1,9329 & 0,9329 & $1{,}42\times10^6$ \\
    11 & 372 & 847,6 & 2,2683 & 1,2683 & $1{,}18\times10^6$ \\
    12 & 326 & 800,7 & 1,9878 & 0,9878 & $1{,}36\times10^6$ \\
    \midrule
    \multicolumn{5}{l}{Среднее значение} & $2{,}17\times10^6$ \\
    \multicolumn{5}{l}{Среднеквадратическое отклонение} & $\pm1{,}04\times10^6$ \\
    \multicolumn{5}{l}{МНК (подгонка)} & $1{,}60\times10^6$ \\
    \bottomrule
  \end{tabular}
\end{table}

\paragraph{Определение среднего $G$ методом наименьших квадратов.}

Поскольку при малых превышениях ($A-1 \ll 1$) погрешность вычисления $G$
существенно возрастает, наилучшую оценку параметра $G$ даёт МНК-подгонка
экспериментальных данных функцией~\eqref{eq:fR}. Получено:
\begin{equation}
  \boxed{G = (1{,}60 \pm 0{,}15)\times 10^6.}
  \label{eq:G_result}
\end{equation}
Среднеарифметическое значение: $\langle G \rangle = (2{,}17 \pm 1{,}04)\times 10^6$.

На рисунке~\ref{fig:frel} показана зависимость $f_R(A)$ с теоретической
кривой при $G = 1{,}60\times10^6$ из МНК-подгонки.

\begin{figure}[H]
  \centering
%   \includegraphics[width=0.85\textwidth]{plot1_frel_vs_A.pdf}
  \begin{tikzpicture}
    \begin{axis}[
        width=12cm,
        height=8cm,
        xlabel={$A$},
        ylabel={$f_R$, кГц},
        grid=both,
        legend pos=north west
    ]

    % Эксперимент
    \addplot[
        only marks,
    ] coordinates {
    (1.10,445.3)
    (1.14,500)
    (1.20,550.7)
    (1.30,613.2)
    (1.42,648.4)
    (1.55,683.5)
    (1.66,699.2)
    (1.75,738.2)
    (1.87,765.6)
    (1.93,796.8)
    (2.27,847.6)
    (1.99,800.7)
    };
    \addlegendentry{Эксперимент}

    % Теория (примерная кривая)
    \addplot[
        domain=1:2.3,
        samples=100,
    ]
    {300*sqrt(x-1)};
    \addlegendentry{Теория}

    \end{axis}
  \end{tikzpicture}
  \caption{Зависимость релаксационной частоты от параметра накачки.
  Точки~--- эксперимент; кривая~--- теория по формуле~\eqref{eq:fR}
  при $G = 1{,}60\times10^6$ (МНК).}
  \label{fig:frel}
\end{figure}

Форма зависимости ($f_R \propto \sqrt{A-1}$) хорошо воспроизводится
теорией; расхождение при малых $A$ объясняется повышенной чувствительностью
$G$ к точности измерения вблизи порога.

\subsection{Задание~2: Ватт-ваттная характеристика лазера}

Снята зависимость мощности лазерного излучения от мощности накачки.
Из всех значений вычтен фоновый сигнал $P_\text{фон} = 0{,}14$~мВт
(измеренный при $P = 0$). Скорректированные данные приведены в
таблице~\ref{tab:task2}.

\begin{table}[H]
  \centering
  \caption{Ватт-ваттная характеристика Nd:YAG лазера (фон вычтен)}
  \label{tab:task2}
  \begin{tabular}{ccccc}
    \toprule
    $P_\text{нак}$, мВт & $P_\text{лаз}^\text{изм}$, мВт &
    $P_\text{лаз}^\text{скорр}$, мВт & $A$ \\
    \midrule
    149 & 0,15 & 0,01  & 0,909 \\
    155 & 0,30 & 0,16  & 0,945 \\
    159 & 0,45 & 0,31  & 0,970 \\
    166 & 0,60 & 0,46  & 1,012 \\
    175 & 0,77 & 0,63  & 1,067 \\
    182 & 0,98 & 0,84  & 1,110 \\
    201 & 1,50 & 1,36  & 1,226 \\
    211 & 2,00 & 1,86  & 1,287 \\
    221 & 2,60 & 2,46  & 1,348 \\
    232 & 3,50 & 3,36  & 1,415 \\
    238 & 4,55 & 4,41  & 1,451 \\
    251 & 5,90 & 5,76  & 1,530 \\
    260 & 6,80 & 6,66  & 1,585 \\
    270 & 7,60 & 7,46  & 1,646 \\
    281 & 8,10 & 7,96  & 1,713 \\
    294 & 8,50 & 8,36  & 1,793 \\
    302 & 8,90 & 8,76  & 1,841 \\
    310 & 9,30 & 9,16  & 1,890 \\
    313 & 9,50 & 9,36  & 1,909 \\
    315 & 9,75 & 9,61  & 1,921 \\
    319 & 10,05 & 9,91 & 1,945 \\
    326 & 10,70 & 10,56 & 1,988 \\
    \bottomrule
  \end{tabular}
\end{table}

По теории~\eqref{eq:stationary} мощность лазера пропорциональна $(A-1)$:
$P_\text{лаз} = k(A-1)$. МНК-подгонка (через ноль) по данным выше порога даёт
наклон:
\begin{equation}
  k = 10{,}42~\text{мВт}.
  \label{eq:k}
\end{equation}

На рисунке~\ref{fig:WW} показана ватт-ваттная характеристика.

\begin{figure}[H]
  \centering
%   \includegraphics[width=0.85\textwidth]{plot2_Plaser_vs_A.pdf}
  \begin{tikzpicture}
    \begin{axis}[
        width=12cm,
        height=8cm,
        xlabel={$A$},
        ylabel={$P_{\text{лаз}}$, мВт},
        grid=both,
        legend pos=north west
    ]

    % Эксперимент (с вычитанием 0.14)
    \addplot[
        only marks,
    ] coordinates {
    (0.91,0.01)
    (0.95,0.16)
    (0.97,0.31)
    (1.01,0.46)
    (1.07,0.63)
    (1.11,0.84)
    (1.23,1.36)
    (1.29,1.86)
    (1.35,2.46)
    (1.41,3.36)
    (1.45,4.41)
    (1.53,5.76)
    (1.59,6.66)
    (1.65,7.46)
    (1.71,7.96)
    (1.79,8.36)
    (1.84,8.76)
    (1.89,9.16)
    (1.91,9.36)
    (1.92,9.61)
    (1.95,9.91)
    (1.99,10.56)
    };
    \addlegendentry{Эксперимент}

    % Теория (линейная зависимость)
    \addplot[
        domain=1:2,
        samples=100,
    ]
    {10*(x-1)};
    \addlegendentry{Теория}

    \end{axis}
  \end{tikzpicture}
  \caption{Ватт-ваттная характеристика Nd:YAG лазера.
  Точки~--- эксперимент (красные~--- выше порога, серые «$\times$»~--- ниже порога);
  прямая~--- теоретическая зависимость $P_\text{лаз} = k(A-1)$.}
  \label{fig:WW}
\end{figure}

Линейный характер ватт-ваттной характеристики выше порога ($A>1$)
хорошо подтверждается экспериментом. Небольшое отклонение вблизи порога
($1 < A < 1{,}2$) связано с тем, что при сильном превышении порога
число лазерных мод ещё мало, а шумы относительно велики.

%%%%%%%%%%%%%%%%%%%%%%%%%%%%%%%%%%%%%%%%%%%%%
\section{Полоса резонатора и сравнение с теорией}

\subsection{Экспериментальная оценка}

По среднему значению $G$ (МНК):
\begin{align}
  T_c &= \frac{T_1}{G} = \frac{0{,}23\times10^{-3}}{1{,}60\times10^6}
        \approx 0{,}144~\text{нс}, \label{eq:Tc}\\
  \delta f_\text{эксп} &= \frac{1}{2\pi T_c} \approx \frac{G}{2\pi T_1}
        \approx 1{,}1~\text{ГГц}. \label{eq:deltaf_exp}
\end{align}

\subsection{Теоретическая оценка}

Для резонатора длиной $L = 6$~см, входное зеркало $R_1 = 1$ (дихроичное),
выходное зеркало $R_2 = 0{,}99$:
\begin{equation}
  \delta f_\text{теор} = -\frac{c\ln\sqrt{R_1 R_2}}{2\pi L}
    = -\frac{3\times10^8 \cdot \ln\sqrt{0{,}99}}{2\pi \cdot 0{,}06}
    \approx 4{,}0~\text{МГц}.
  \label{eq:deltaf_teor_val}
\end{equation}

\subsection{Сравнение}

Экспериментальная полоса резонатора $\delta f_\text{эксп} \approx 1{,}1$~ГГц
значительно превышает теоретическую $\delta f_\text{теор} \approx 4$~МГц.
Это расхождение обусловлено тем, что в реальном резонаторе суммарные потери
складываются не только из потерь на выходном зеркале, но и из:
\begin{itemize}
  \item потерь на рассеяние и поглощение в кристалле Nd:YAG;
  \item дифракционных потерь;
  \item потерь в оптических элементах (фильтр, линзы);
  \item нежелательных отражений и юстировочных погрешностей.
\end{itemize}
Таким образом, полные потери в данном резонаторе значительно выше, чем
обусловленные одним только выходным зеркалом, что и приводит к существенно
большей экспериментальной полосе.

%%%%%%%%%%%%%%%%%%%%%%%%%%%%%%%%%%%%%%%%%%%%%
\section{Оценка погрешностей}

\subsection{Погрешность параметра накачки $A$}

Инструментальная погрешность измерения мощности $\Delta P = 2$~мВт:
\begin{equation}
  \Delta A = \frac{\Delta P}{P_\text{пор}} = \frac{2}{164} \approx 0{,}012.
\end{equation}

\subsection{Погрешность безразмерной частоты $\Omega_R$}

При погрешности частотомера $\Delta f = 5$~кГц:
\begin{equation}
  \Delta\Omega_R = 2\pi \cdot \Delta f \cdot T_1 = 2\pi \cdot 5\times10^3 \cdot 0{,}23\times10^{-3}
  \approx 7{,}2\times10^{-3}.
\end{equation}

\subsection{Погрешность $G$}

Поточечно из формулы $G = \Omega_R^2/(A-1)$, по формуле распространения ошибок:
\begin{equation}
  \frac{\Delta G}{G} = \sqrt{\left(\frac{2\Delta\Omega_R}{\Omega_R}\right)^2
    + \left(\frac{\Delta A}{A-1}\right)^2}.
\end{equation}
Для средних значений $A \approx 1{,}5$, $\Omega_R \approx 900$:
\begin{equation}
  \frac{\Delta G}{G} \approx \sqrt{(0{,}016)^2 + (0{,}024)^2} \approx 0{,}029,
  \quad \Delta G \approx 0{,}029 \cdot G \approx 5\times10^4.
\end{equation}

Основным источником неопределённости является большой разброс $G_i$,
вычисленных поточечно (особенно вблизи порога), поэтому МНК-оценка
$G_\text{МНК} = (1{,}60 \pm 0{,}15)\times10^6$ является более надёжной.

%%%%%%%%%%%%%%%%%%%%%%%%%%%%%%%%%%%%%%%%%%%%%
\section{Сводная таблица результатов}

\begin{table}[H]
  \centering
  \caption{Итоговые результаты лабораторной работы}
  \label{tab:results}
  \begin{tabular}{lll}
    \toprule
    Величина & Эксперимент & Теория/норма \\
    \midrule
    Пороговая мощность накачки $P_\text{пор}$ & $164$~мВт & --- \\
    Параметр $G$ (МНК) & $(1{,}60 \pm 0{,}15)\times10^6$ & $\sim 10^3$--$10^5$ \\
    Параметр $G$ (среднеарифм.) & $(2{,}17 \pm 1{,}04)\times10^6$ & --- \\
    Время жизни фотона $T_c$ & $0{,}14$~нс & $\sim 40$~нс \\
    Полоса резонатора $\delta f_\text{эксп}$ & $\approx 1{,}1$~ГГц & --- \\
    Полоса резонатора $\delta f_\text{теор}$ & --- & $\approx 4$~МГц \\
    Наклон ВВХ $k$ & $10{,}42$~мВт & --- \\
    \bottomrule
  \end{tabular}
\end{table}

%%%%%%%%%%%%%%%%%%%%%%%%%%%%%%%%%%%%%%%%%%%%%
\section{Выводы}

\begin{enumerate}
  \item Экспериментально подтверждена квадратично-корневая зависимость
    релаксационной частоты от параметра накачки:
    $f_R \propto \sqrt{A-1}$, что соответствует теории ТСД.

  \item Линейный характер ватт-ваттной характеристики ($P_\text{лаз} \propto A-1$)
    выше порога генерации подтверждается экспериментом.

  \item Из данных по релаксационным колебаниям определён параметр
    $G = (1{,}60 \pm 0{,}15)\times10^6$ и экспериментальная полоса
    резонатора $\delta f_\text{эксп} \approx 1{,}1$~ГГц.

  \item Экспериментальная полоса резонатора примерно в 275 раз превышает
    теоретическую оценку по потерям на зеркалах ($\delta f_\text{теор} \approx 4$~МГц).
    Это свидетельствует о значительных дополнительных потерях в реальном
    резонаторе (рассеяние в кристалле, дифракционные потери и т.\,д.).

  \item Установленный порог генерации $P_\text{пор} = 164$~мВт согласуется
    с ватт-ваттными данными: при мощности выше 164~мВт наблюдается
    устойчивая лазерная генерация.
\end{enumerate}

%%%%%%%%%%%%%%%%%%%%%%%%%%%%%%%%%%%%%%%%%%%%%
\begin{thebibliography}{9}
  \bibitem{zveltо} Звелто О. \textit{Принципы лазеров}. М.: Мир, 1990.
  \bibitem{khanin} Ханин Я.\,И. \textit{Динамика квантовых генераторов}. М.: Сов. Радио, 1975.
  \bibitem{tang1963} Tang C.\,L., Statz H. and deMars G. Spectral output and spiking
    behaviour of solid-state laser // \textit{J.\,Appl.\,Phys.}\ 1963. V.\,34, №8. P.\,2289.
  \bibitem{rabinovich} Рабинович М.\,И., Трубецков Д.\,И. \textit{Введение в теорию
    колебаний и волн}. М.: Наука, 1984.
  \bibitem{yariv} Ярив А. \textit{Введение в оптическую электронику}. М.: Высшая школа, 1983.
\end{thebibliography}

%%%%%%%%%%%%%%%%%%%%%%%%%%%%%%%%%%%%%%%%%%%%%
\appendix
\section{Протокол измерений}

\subsection*{Задание~1: Релаксационная частота}

\begin{table}[H]
  \centering
  \begin{tabular}{ccc}
    \toprule
    № & Мощность накачки $P$, мВт & Релаксационная частота $f_R$, кГц \\
    \midrule
    1  & 180 & 445,3 \\
    2  & 187 & 500,0 \\
    3  & 196 & 550,7 \\
    4  & 214 & 613,2 \\
    5  & 233 & 648,4 \\
    6  & 254 & 683,5 \\
    7  & 272 & 699,2 \\
    8  & 287 & 738,2 \\
    9  & 307 & 765,6 \\
    10 & 317 & 796,8 \\
    11 & 372 & 847,6 \\
    12 & 326 & 800,7 \\
    \bottomrule
  \end{tabular}
  \caption{Протокол измерений. Задание~1.}
\end{table}

\subsection*{Задание~2: Ватт-ваттная характеристика}
Фоновый сигнал: $P_\text{фон} = 0{,}14$~мВт (вычтен из всех значений).

\begin{table}[H]
  \centering
  \begin{tabular}{cc}
    \toprule
    Мощность накачки $P_\text{нак}$, мВт & Мощность лазера $P_\text{лаз}^\text{изм}$, мВт \\
    \midrule
    149 & 0,15 \\
    155 & 0,30 \\
    159 & 0,45 \\
    166 & 0,60 \\
    175 & 0,77 \\
    182 & 0,98 \\
    201 & 1,50 \\
    211 & 2,00 \\
    221 & 2,60 \\
    232 & 3,50 \\
    238 & 4,55 \\
    251 & 5,90 \\
    260 & 6,80 \\
    270 & 7,60 \\
    281 & 8,10 \\
    294 & 8,50 \\
    302 & 8,90 \\
    310 & 9,30 \\
    313 & 9,50 \\
    315 & 9,75 \\
    319 & 10,05 \\
    326 & 10,70 \\
    \midrule
    0   & 0,14 (фон) \\
    \bottomrule
  \end{tabular}
  \caption{Протокол измерений. Задание~2.}
\end{table}

\end{document}
